\section*{Chapitre \ref{ch:b3:adaptive-immunity} --- Immunité adaptative et vaccination}

\begin{solution}{exo:b3:adaptive-immunity:1}
Un \omterm{def:b3:adaptive-immunity:lymphocytes}{lymphocyte}, une spécificité (Nossal et Lederberg~: chaque cellule ne
sécrétait d'\omterm{def:b3:adaptive-immunity:antibody}{anticorps} que contre un seul \omterm{def:b3:adaptive-immunity:lymphocytes}{antigène})~; le récepteur est
fabriqué avant la rencontre de l'\omterm{def:b3:adaptive-immunity:lymphocytes}{antigène} (Tonegawa~: le gène est
réarrangé dans le \omterm{def:b3:adaptive-immunity:lymphocytes}{lymphocyte}~B, non en réponse à l'\omterm{def:b3:adaptive-immunity:lymphocytes}{antigène})~;
l'\omterm{def:b3:adaptive-immunity:lymphocytes}{antigène} sélectionne et amplifie les clones qui conviennent en
cellules effectrices et mémoires~; les clones autoréactifs sont
éliminés pendant le développement.
\end{solution}

\begin{solution}{exo:b3:adaptive-immunity:2}
Deux chaînes lourdes et deux légères, chacune avec un domaine variable
à la pointe et des domaines constants derrière~; deux sites de liaison
identiques formés par six boucles hypervariables~; une tige Fc que
lisent les autres cellules et le \omterm{def:b3:innate-immunity:complement}{complément}. IgM~: premier \omterm{def:b3:adaptive-immunity:antibody}{anticorps},
pentamère, activation du \omterm{def:b3:innate-immunity:complement}{complément}. IgG~: principal \omterm{def:b3:adaptive-immunity:antibody}{anticorps} du sang
et des tissus, \omterm{def:b3:innate-immunity:complement}{opsonisation}, neutralisation, traverse le placenta.
IgA~: dimère sécrété sur les surfaces muqueuses et dans le lait. IgE~:
liée aux mastocytes, contre les parasites, médiatrice de l'allergie.
IgD~: sur les \omterm{def:b3:adaptive-immunity:lymphocytes}{lymphocytes}~B naïfs, de fonction connue mineure.
\end{solution}

\begin{solution}{exo:b3:adaptive-immunity:3}
Classe~I~: toutes les cellules nucléées~; peptides de protéines
cytosoliques découpées par le protéasome~; 8 à 10 résidus~; lus par les
\omterm{def:b3:adaptive-immunity:lymphocytes}{lymphocytes}~T cytotoxiques CD8. Classe~II~: \omterm{def:b3:innate-immunity:innate}{cellules dendritiques},
\omterm{def:b3:innate-immunity:innate}{macrophages}, \omterm{def:b3:adaptive-immunity:lymphocytes}{lymphocytes}~B~; peptides de protéines internalisées et
digérées dans les \omterm{def:b3:membrane-traffic:endocytosis}{endosomes}~; 13 à 25 résidus~; lus par les
\omterm{def:b3:adaptive-immunity:lymphocytes}{lymphocytes}~T auxiliaires CD4.
\end{solution}

\begin{solution}{exo:b3:adaptive-immunity:4}
$R_{0}$ est le nombre moyen de cas qu'un cas produit dans une population
entièrement susceptible~; le seuil est $p_{c} = 1 - 1/R_{0}$. Oreillons
\qty{80}{\%}~; coqueluche \qty{92}{\%}~; grippe de 1918 \qty{60}{\%}.
\end{solution}

\begin{solution}{exo:b3:adaptive-immunity:5}
$50\times 20\times 5 = 5000$ chaînes lourdes, $60\times 4 = 240$
légères~: $1.2\times 10^{6}$ combinaisons~; avec la \omterm{def:b3:adaptive-immunity:antibody}{diversité
jonctionnelle} $1.2\times
10^{10}$ --- dix fois le nombre de \omterm{def:b3:adaptive-immunity:lymphocytes}{lymphocytes}~B, de sorte que l'animal
exprime au plus un dixième de ce qu'il pourrait faire, et que presque
chaque cellule porte un récepteur unique.
\end{solution}

\begin{solution}{exo:b3:adaptive-immunity:6}
$1000\times 10^{-3} = 1$ mutation par cellule et par division~; $20$ en
vingt divisions. Parmi $10^{5}$ cellules, $2\times 10^{6}$ événements de
mutation tombent sur $3000$ substitutions simples possibles~: chacune
est échantillonnée des centaines de fois, et les combinaisons doubles et
triples par-dessus le marché --- le \omterm{def:b3:adaptive-immunity:germinal-centre}{centre germinatif} explore le
voisinage du récepteur de façon exhaustive.
\end{solution}

\begin{solution}{exo:b3:adaptive-immunity:7}
Les \omterm{def:b3:adaptive-immunity:lymphocytes}{lymphocytes}~T sont sélectionnés dans le \omterm{def:b3:adaptive-immunity:tolerance}{thymus} pour reconnaître des
peptides sur les allèles du CMH de l'hôte lui-même~; les allèles d'une
autre personne ont des sillons de forme différente qui retiennent
d'autres peptides dans d'autres orientations, de sorte que le même
peptide viral n'est pas vu. Les molécules du CMH étrangères d'un
greffon sont elles-mêmes reconnues par une grande partie des
\omterm{def:b3:adaptive-immunity:lymphocytes}{lymphocytes}~T comme «~le CMH du soi portant un peptide étrange~», d'où
le rejet. Le polymorphisme persiste parce qu'une population à nombreux
allèles présente un échantillon plus large des peptides de tout
pathogène et qu'aucun pathogène ne peut échapper à tout le monde~; les
hétérozygotes présentent deux fois plus et sont favorisés.
\end{solution}

\begin{solution}{exo:b3:adaptive-immunity:8}
(a) Pas de \omterm{def:b3:adaptive-immunity:antibody}{recombinaison V(D)J}~: aucun \omterm{def:b3:adaptive-immunity:lymphocytes}{lymphocyte B} ni T, déficit
immunitaire combiné sévère. (b) Pas d'AIRE~: les cellules thymiques
n'exposent pas les \omterm{def:b3:adaptive-immunity:lymphocytes}{antigènes} des tissus, les \omterm{def:b3:adaptive-immunity:lymphocytes}{lymphocytes}~T autoréactifs
s'échappent, \omterm{def:b3:adaptive-immunity:tolerance}{auto-immunité} contre de nombreux organes. (c) Pas de
\omterm{def:b3:adaptive-immunity:tolerance}{FoxP3}~: aucun \omterm{def:b3:adaptive-immunity:lymphocytes}{lymphocyte}~T régulateur, \omterm{def:b3:adaptive-immunity:tolerance}{auto-immunité} et allergie
précoces et mortelles. (d) Pas d'AID~: ni \omterm{def:b3:adaptive-immunity:antibody}{commutation de classe} ni
hypermutation --- IgM abondante, presque pas d'IgG ni d'IgA, \omterm{def:b3:adaptive-immunity:antibody}{anticorps}
de faible affinité, infections à répétition. (e) Pas de \omterm{def:b3:adaptive-immunity:lymphocytes}{lymphocytes}~T
CD4~: aucune aide pour les \omterm{def:b3:adaptive-immunity:lymphocytes}{lymphocytes}~B ni pour les \omterm{def:b3:innate-immunity:innate}{macrophages},
réponses \omterm{def:b3:adaptive-immunity:antibody}{anticorps} médiocres et sensibilité aux infections
opportunistes --- le déficit immunitaire du sida.
\end{solution}

\begin{solution}{exo:b3:adaptive-immunity:9}
$p_{c} = 1 - 1/6 = 0.83$~; couverture $0.83/0.9 = 93\,\%$. À \qty{80}{\%}
de couverture la fraction immune est $0.72$, $R_{\text{eff}} = 6\times 0.28
= 1.7$. Parmi les susceptibles l'épidémie se propage comme si $R_{0}$
valait $1.7$ (les immuns ne font que diluer les contacts), et la
relation de taille finale $s_{\infty} - \ln s_{\infty}/1.7 = 1$ donne
$s_{\infty}
\approx 0.31$~: environ \qty{69}{\%} des susceptibles, \qty{19}{\%} de
la population, sont infectés.
\end{solution}

\begin{solution}{exo:b3:adaptive-immunity:10}
À partir de $\mathrm{d}s/\mathrm{d}t = -\beta si$ et $\mathrm{d}i/\mathrm{d}t
= \beta si - \gamma i$, $\mathrm{d}i/\mathrm{d}s = -1 + \gamma/(\beta s)$,
donc $i + s - (\ln s)/R_{0}$ est conservé~; avec $s = 1$, $i = 0$ au
départ et $i = 0$, $s = s_{\infty}$ à la fin, $s_{\infty} - \ln
s_{\infty}/R_{0} = 1$. Solutions~: $R_{0} = 1.5$, $s_{\infty} = 0.42$
(\qty{58}{\%} d'infectés)~; $2$, $0.20$ (\qty{80}{\%})~; $3$, $0.06$
(\qty{94}{\%})~; $5$, $0.007$ (\qty{99.3}{\%}). Pas tout le monde~: à
mesure que les susceptibles s'épuisent le nombre de reproduction
effectif tombe sous un alors qu'il en reste, et l'épidémie s'éteint
avant de les atteindre.
\end{solution}

\begin{solution}{exo:b3:adaptive-immunity:11}
$1.5^{n} = 1000$~: $n = \ln 1000/\ln 1.5 \approx 17$ améliorations
successives. Avec une mutation par cellule et par division et une sur
cent qui améliore, il faut cribler au moins cent cellules par tour pour
en trouver une~; en pratique des milliers, puisque la moitié des
mutations détruit la liaison et que la sélection est bruitée. Dix-sept
tours de deux divisions chacun prennent environ neuf jours --- les deux
à trois semaines observées, compte tenu des inefficacités. Les cycles
séparent la mutation (zone sombre) du test (zone claire), de sorte que
chaque variant est jugé face à l'\omterm{def:b3:adaptive-immunity:lymphocytes}{antigène} avant qu'on le laisse se
multiplier, exactement comme un algorithme génétique alterne variation
et sélection.
\end{solution}

\begin{solution}{exo:b3:adaptive-immunity:12}
Les femmes~: plusieurs gènes immunitaires sont sur le chromosome X,
certains (TLR7) échappent à l'inactivation et s'expriment à partir des
deux copies, et les œstrogènes améliorent la survie des \omterm{def:b3:adaptive-immunity:lymphocytes}{lymphocytes}~B
--- test~: les hommes porteurs d'un X supplémentaire (Klinefelter)
devraient avoir un risque de lupus proche de celui des femmes, et c'est
le cas. \omterm{def:b3:adaptive-immunity:mhc}{HLA}~: les allèles associés présentent bien les peptides du soi
en cause, ou ne les présentent pas dans le \omterm{def:b3:adaptive-immunity:tolerance}{thymus} de sorte que les
clones ne sont pas éliminés --- test~: des souris \omterm{def:b3:genetic-engineering:transgenic}{transgéniques} pour
l'allèle humain développent la maladie quand on leur donne l'\omterm{def:b3:adaptive-immunity:lymphocytes}{antigène}.
La hausse~: une moindre exposition microbienne précoce et des
\omterm{def:b3:microbiomes:symbiosis}{microbiotes} altérés (\omterm{def:b3:bacteriology:antibiotics}{antibiotiques}, régime, naissance par césarienne)
donnent moins de \omterm{def:b3:adaptive-immunity:lymphocytes}{lymphocytes}~T régulateurs et plus d'\omterm{def:b3:innate-immunity:inflammation}{inflammation} ---
tests~: les souris prédisposées au diabète élevées sans germes ou sous
\omterm{def:b3:bacteriology:antibiotics}{antibiotiques} en développent davantage, et les enfants d'immigrés
acquièrent l'incidence du pays d'accueil en une génération.
\end{solution}

\begin{solution}{pb:b3:adaptive-immunity:1}
\textbf{1.} $6000\times 320 = 1.9\times 10^{6}$~; $\times 3\times 10^{4}
= 5.8\times 10^{10}$.
\textbf{2.} $10^{11}$ cellules face à $5.8\times 10^{10}$ possibles~: un
corps pourrait contenir chaque récepteur environ deux fois, mais les
clones sont inégaux et la plupart des séquences ne sont jamais
fabriquées, si bien que chaque corps exprime un échantillon vaste mais
incomplet --- et deux personnes ne partagent qu'une petite fraction de
leurs récepteurs.
\textbf{3.} $10^{11}/10^{5} = 10^{6}$ cellules~; dans un ganglion de
$10^{8}$, environ $1000$.
\textbf{4.} $10^{9}/10^{5} = 10^{4}$ précurseurs --- cent fois moins~; le
nouveau-né a aussi des follicules et des \omterm{def:b3:innate-immunity:innate}{cellules dendritiques}
immatures, aucune mémoire et une aide limitée des \omterm{def:b3:adaptive-immunity:lymphocytes}{lymphocytes}~T, de
sorte que ses réponses sont lentes et faibles, et l'IgG maternelle
comble l'intervalle.
\textbf{5.} Elle enjambe les jonctions V--D et D--J, où des nucléotides
sont retirés et ajoutés au hasard~: sa séquence n'est codée nulle part
dans le génome. Assise entre les deux chaînes au centre du site de
liaison, c'est la boucle le plus souvent au contact de l'épitope.
\textbf{6.} Avant l'\omterm{def:b3:adaptive-immunity:lymphocytes}{antigène}, la diversité doit être large pour que
quelque chose lie ce qui viendra~; après l'\omterm{def:b3:adaptive-immunity:lymphocytes}{antigène}, l'effort est mieux
employé à varier les quelques récepteurs dont on sait déjà qu'ils lient.
Une recherche grossière puis une recherche fine valent bien mieux que
l'une ou l'autre seule.
\textbf{7.} Sept jours font $21$ divisions~: $100\times 2^{21} = 2.1\times
10^{8}$ cellules~; \qty{30}{\%}, $6.3\times 10^{7}$ \omterm{def:b3:adaptive-immunity:response}{plasmocytes}.
\textbf{8.} $6.3\times 10^{7}\times 2000\times \num{86400} = 1.1\times
10^{16}$ molécules par jour~; $\times 1.5\times 10^{5}\times 1.66\times
10^{-24}$ g $= \qty{2.7}{mg}$~; dans \qty{3}{L}, \qty{0.9}{mg/L} par jour.
\textbf{9.} \qty{27}{mg} dans \qty{3}{L}~: \qty{9}{mg/L} $= \qty{0.009}{g/L}$,
un millième de l'IgG totale --- l'\omterm{def:b3:adaptive-immunity:antibody}{anticorps} dirigé contre un seul
\omterm{def:b3:adaptive-immunity:lymphocytes}{antigène} est une petite fraction de l'ensemble.
\textbf{10.} Un facteur cent fait $\log_{2}100 = 6.6$ demi-vies, environ
$140$ jours. Les années d'\omterm{def:b3:adaptive-immunity:antibody}{anticorps} viennent des \omterm{def:b3:adaptive-immunity:response}{plasmocytes} de longue
vie de la moelle, qui sécrètent en continu, et des \omterm{def:b3:adaptive-immunity:response}{cellules mémoires}
restimulées à la réexposition.
\textbf{11.} Neuf divisions~: $10^{5}\times 512 = 5\times 10^{7}$ cellules
au jour 3 --- mille fois les $5\times 10^{4}$ de la primaire au jour 3,
et le quart de ce que la primaire n'atteignait qu'au jour 7.
\textbf{12.} Le VDJ réarrangé est voisin de C$_{\mu}$, de sorte que le
premier transcrit est une IgM~; commuter vers C$_{\gamma}$ demande
l'aide des \omterm{def:b3:adaptive-immunity:lymphocytes}{lymphocytes}~T et AID, ce qui prend des jours. Les \omterm{def:b3:adaptive-immunity:response}{cellules
mémoires} ont déjà commuté, donc leur descendance fait d'emblée de l'IgG.
\textbf{13.} $1000\times 10^{-3} = 1$ par division~; $20$ sur vingt.
\textbf{14.} Une sur cent~; $10^{4}\times 0.01 = 100$ cellules améliorées
par division.
\textbf{15.} $1.5^{n} = 1000$~: $n \approx 17$~; $17\times \qty{12}{h}
\approx 8.5$ jours.
\textbf{16.} La moitié des filles perd la liaison~: sans sélection, $0.5^{20}
\approx 10^{-6}$ de la lignée lierait encore après vingt divisions. La
zone claire doit retirer les perdantes à chaque cycle pour que seuls des
récepteurs fonctionnels et améliorés poursuivent.
\textbf{17.} La première dose laisse des \omterm{def:b3:adaptive-immunity:lymphocytes}{lymphocytes}~B mémoires
d'affinité relevée et de classe commutée~; la seconde dose ensemence
avec eux de nouveaux \omterm{def:b3:adaptive-immunity:germinal-centre}{centres germinatifs}, et l'hypermutation et la
sélection reprennent d'un point de départ plus haut, si bien que les
\omterm{def:b3:adaptive-immunity:antibody}{anticorps} qui en sortent sont meilleurs encore.
\textbf{18.} Certaines personnes infectées finissent par fabriquer des
\omterm{def:b3:adaptive-immunity:antibody}{anticorps} contre des sites conservés que le \omterm{def:b3:virology:virus}{virus} ne peut changer
(\omterm{def:b3:adaptive-immunity:antibody}{anticorps} largement neutralisants), après des années de maturation. Le
\omterm{def:b3:adaptive-immunity:germinal-centre}{centre germinatif} peut être piloté~: une série d'immunogènes qui
engagerait d'abord les bons précurseurs germinaux puis récompenserait la
liaison au site conservé pourrait guider la maturation sur ce chemin en
quelques mois plutôt qu'en années.
\textbf{19.} $p_{c} = 1 - 1/15 = 93.3\,\%$. Une dose~: $0.933/0.93 =
100.4\,\%$ --- hors d'atteinte~; deux doses~: $0.933/0.97 = 96.2\,\%$.
\textbf{20.} Immuns $0.90\times 0.97 = 0.873$~; $R_{\text{eff}} = 15\times
0.127 = 1.9$~: la rougeole circule.
\textbf{21.} $60\,000\times (0.10 + 0.90\times 0.03) = 7600$ susceptibles
par cohorte de naissance~; après cinq ans $38\,000$ --- un réservoir où
un cas importé trouve dans ses contacts assez de susceptibles pour
entretenir des chaînes.
\textbf{22.} Parmi les susceptibles, $s_{\infty} - \ln s_{\infty}/1.9 =
1$ donne $s_{\infty} \approx 0.23$~: environ \qty{77}{\%} d'entre eux,
soit quelque $29\,000$ cas.
\textbf{23.} Immuns $0.95\times 0.97 = 0.92$~; $R_{\text{eff}} = 15\times
0.08 = 1.2$ --- toujours au-dessus de un~; le seuil demande \qty{96}{\%}
avec deux doses, et la rougeole est la maladie qui punit les derniers
pour cent.
\textbf{24.} Les nourrissons ne sont protégés que par l'immunité de leur
entourage (et par l'\omterm{def:b3:adaptive-immunity:antibody}{anticorps} maternel pendant quelques mois)~: si les
chaînes de transmission ne peuvent se former, aucun cas ne les atteint.
À \qty{85}{\%} de couverture la fraction immune est $0.82$ et
$R_{\text{eff}} = 2.6$~: les flambées courent, et les nourrissons, non
vaccinés et les plus vulnérables, sont infectés.
\textbf{25.} Répertoire d'environ $6\times 10^{10}$~; \qty{2.7}{mg} d'IgG
par jour venant des \omterm{def:b3:adaptive-immunity:response}{plasmocytes} d'une seule réponse~; couverture à deux
doses de \qty{96}{\%}.
\end{solution}
